\documentclass{article}
\usepackage{ctex}
\usepackage{amssymb}
\usepackage{amsmath}
\title{线性回归}
\author{QwanxingxuA}


\begin{document}
\maketitle
线性回归是很简单的模型，也是比较好实现的。李航老师详细介绍
了线性回归的方法，但是策略层面对我这种初学者不友好，尤其是
概率模型，以下说说我的思路。

\section{模型}
对于样本空间$D=\{(x_1,y_1),(x_2,y_2),\dots,(x_n,y_n)\}$,输入
空间$\aleph\subseteq \mathbb{R}^m $,输出空间$\mathcal{Y}\subseteq \mathbb{R}$。基于对数据的初步观察
，缩小假设空间到\[y=\boldsymbol{w}\cdot \boldsymbol{x}+b\]

构建模型的本质就是基于对数据的初步观察缩小假设空间。

\section{策略}
策略是基于已经缩小的假设空间找到最优$f$，通过经验误差或者说损失
误差来找到最优，同时尽量让风险误差或者泛化误差小。
\subsection{直观非概率}
从直观上分析经验误差就是预测值$\widehat{y}$与真实值$y$的
距离。显然可以得到损失函数\[L(\boldsymbol{w})=\frac{1}{2n}\sum_{i=1}^{n}
(y_i-\widehat{y_i})^2=\frac{1}{2n}\sum_{i=1}^{n}
[y_i-(\boldsymbol{w}\cdot \boldsymbol{x_i}+b)]^2\]
该损失函数优化求解就是线性最小二乘法，见最优化方法：

这里说明一下$X^TX$半正定（这也是李航老师书中题目）：
对任意非零向量$\boldsymbol{u}$有：
\[\boldsymbol{u}^TX^TX\boldsymbol{u}=(X\boldsymbol{u})^TX\boldsymbol{u}\ge 0\]

\subsection{正态分布概率}
这是这一章不好理解的点。实际上，这种策略是假设条件概率密度$p(x,y)$服从正态
分布。不好理解的关键是与直观距离相比，这里的目标是使得输入
$x$的真实值靠近正态分布的最高点。
\[E(y|\boldsymbol{x})=E(\boldsymbol{w}\cdot \boldsymbol{x}+b|\boldsymbol{x})
=\boldsymbol{w}\cdot \boldsymbol{x}+b=\widehat{y}\]
从均值来看，也可以说是优化参数使得输入$x$的预测值$\widehat{y}$
在正态分布概率下靠近真实值$y$。

最大似然估计是为了让观察到的数据出现可能性最大，这天然
适配上述问题。代入此问题就是使真实数据处于均值，这样数据
在真实值附近概率最大。

取似然函数的负数，构成损失函数：
\[L(\boldsymbol{w})=\frac{1}{2}\log(2\pi \sigma^2)+\frac{1}{2
\sigma ^2}\sum_{i=1}^{n}[y_i-(\boldsymbol{w}\cdot \boldsymbol{x_i}+b)]^2)]\]
损失函数是用来服务策略的，是个优化问题，可以忽视常数和一些常量，自然可以得到：
\[L(\boldsymbol{w})=\frac{1}{2n}\sum_{i=1}^{n}
[y_i-(\boldsymbol{w}\cdot \boldsymbol{x_i}+b)]^2\]
这里回到最小二乘法了，李航老师的书就是这么写的，很容易使初学者如我一样忽略
概率，凭借直观感受忽略了概率策略的思想。

\section{结构风险误差}
线性回归模型还有岭回归和Lasso，它们是两种正则化手段，也就是
结构风险误差，实质上还是一种策略。

但是两种损失函数对泛化误差产生了不同影响，这也印证了我在二分类任务
泛化误差上界中说的损失函数与泛化误差存在着联系。

总之，线性回归模型是很简单的模型，其中的概率策略实质上就是
最大似然估计的思想。作者似乎记得哪个题目要求证明风险误差就是最大似然估计来着，
现在的作者实力还不够，慢慢学习理解吧。
\end{document}